\documentclass[10pt,oneside]{book}
%%%% Page Info + Commands %%%%%{

%packages
\usepackage{pgfplots}
\usepackage{mathrsfs}
\pgfplotsset{compat=1.18}
\usepackage{amsmath}
\usepackage{amsfonts}
\usepackage{charter}
\usepackage{amsthm,letltxmacro}
\usepackage{amssymb}
\usepackage{fancyhdr}
\usepackage{float}
\usepackage[most]{tcolorbox}
\usepackage{enumerate}
\usepackage{enumitem}
\usepackage[dvipsnames]{xcolor}
\usepackage{changepage}
\usepackage{mdframed}
\usepackage{graphicx}

% SMALL PAGE COMMAND
\newcommand{\smallpage}{  \setlength \oddsidemargin{.25in}
            \setlength \textwidth{5.5in}
            \setlength \topmargin{0in}
            \setlength \textheight{9in}}


% Commands for sets of numbers
\newcommand{\Z}{\mathbb{Z}}\newcommand{\R}{\mathbb{R}}\newcommand{\N}{\mathbb{N}}\newcommand{\Q}{\mathbb{Q}}\newcommand{\C}{\mathbb{C}}
\newcommand{\real}[1]{\text{Re}\left(#1\right)}
\newcommand{\im}[1]{\text{Im}\left(#1\right)}

% gordie spacing and boxing commands
\newcommand{\gspc}{\vspace{0.13cm}} % 0.5 space
\newcommand{\gline}[0]{\gspc\\} % 1.5 space
\newcommand{\gbline}{\gspc\gspc\\} % double space
\newcommand{\gbox}[1]{\fbox{\parbox{\linewidth}{#1}}} % multiline box command
\newcommand{\gmod}[1]{\,(\text{mod}\,#1)}


\newcommand{\gimage}[3]{
\begin{figure}[H]
    \centering
    \includegraphics[width=#2]{#1}
    \caption{#3}
    \label{fig:#1}
\end{figure}
}

\newcommand{\centerit}[1]{{\begin{center} #1 \end{center}}}
\newcommand{\ul}[1]{\underline{#1}}

\newcommand{\prt}[2]{\frac{\partial#1}{\partial#2}}

%%%%%%%% command for graphics %%%%%%%%%%%%%

%}

% COLOR BOX
\newmdenv[backgroundcolor=gray!8,
  linewidth=0pt, innerleftmargin=0pt, innerrightmargin=0pt,
  skipabove=\baselineskip, skipbelow=\baselineskip
]{coloredadjust}

% PROOF
\LetLtxMacro\oldproof\proof
\let\endoldproof\endproof
\renewenvironment{proof}[1][\proofname]
  {\vspace{0.2cm}\begin{adjustwidth}{2em}{2em}
   \oldproof[#1]}
  {\endoldproof
\end{adjustwidth}}

% PROBLEM
\newenvironment{problem}[1]
  {\vspace{-0.1cm}\begin{coloredadjust}\begin{adjustwidth}{2em}{2em}
   \textit{Problem. }#1}
  {\endoldproof
\end{adjustwidth}\end{coloredadjust}\gspc}

%% DEFINITION
\newtcbtheorem[]{definition}{Definition}{enhanced,
  colback=blue!2, colframe=blue!50!black, coltitle=blue!70!black,
  fonttitle=\bfseries,
  boxrule=0pt, toprule=2pt, bottomrule=1pt, leftrule=1pt, rightrule=1pt, arc=1pt, outer arc=1pt,
  attach title to upper, title={#1},
}{dfn}

%% CRITERION
\newtcbtheorem[]{criterion}{Criterion}{enhanced,
  colback=magenta!3, colframe=magenta!60!black, coltitle=magenta!20!black,
  fonttitle=\bfseries,
  boxrule=4pt, toprule=2pt, bottomrule=1pt, leftrule=1pt, rightrule=1pt, arc=1pt, outer arc=1pt,
  attach title to upper, title={#1},
}{ctn}

%% COROLLARY
\newtcbtheorem[]{corollary}{Corollary}{enhanced,
  colback=orange!3, colframe=orange!80!black, coltitle=orange!50!black,
  fonttitle=\bfseries,
  boxrule=4pt, toprule=1pt, bottomrule=1pt, leftrule=1pt, rightrule=1pt, arc=5pt, outer arc=5pt,
  attach title to upper, title={#1},
}{ctn}

%% THEOREM
\newtcbtheorem[]{theorem}{Theorem}{enhanced,
  colback=green!2, colframe=green!40!black, coltitle=green!20!black,
  fonttitle=\bfseries,
  boxrule=4pt, toprule=3pt, bottomrule=1pt, leftrule=1pt, rightrule=1pt,arc=5pt, outer arc=5pt,
  attach title to upper, title={#1},
}{thm}

%%%% Page Setup %%%%%%%
\smallpage\sloppy
\pagestyle{fancy}
\lhead{\large{\textbf{PS1 - Maxwell's Equations}}}
\setlength{\headheight}{10pt}


% color commands
\newcommand{\cg}[1]{\color{OliveGreen}#1\color{white}}
\newcommand{\cb}[1]{\color{BlueViolet}#1\color{white}}


\newcommand{\cpr}[1]{\left(#1\right)}
%%%%%%%%%%%%% PHYSICS SPECIFIC COMMANDS

\newcommand{\vA}{\vec{\mathbf{A}}}
\newcommand{\vB}{\vec{\mathbf{B}}}
\newcommand{\vC}{\vec{\mathbf{C}}}
\newcommand{\vZ}{\vec{\mathbf{0}}}
\newcommand{\norm}{\vec{\mathbf{n}}}
\newcommand{\pvec}[1]{\left\langle#1\right\rangle}
\newcommand{\ar}{\vec{\mathbf{r}}}
\newcommand{\arhat}{\hat{\mathbf{r}}}
\newcommand{\vv}{\vec{\mathbf{v}}}
\newcommand{\delfull}{\pvec{\prt{}{x}, \prt{}{y},\prt{}{z}}}
\newcommand{\delinput}[3]{\pvec{\prt{}{x}\cpr{#1}, \prt{}{y}\cpr{#2},\prt{}{z}\cpr{#3}}}
\newcommand{\dps}{\displaystyle}
\newcommand{\delmatrix}[3]{\begin{bmatrix}\dps\hat x&\dps\hat y&\dps\hat z\\\dps\prt{}x&\dps\prt{}y&\dps\prt{}z\\\dps#1&\dps#2&\dps#3\end{bmatrix}}

%%%%%%%%%%%%%%%%%%%%%%%%%%%%
%%%%%%%%%%%%%%%%%%%%%%%%%%%%%%
%%%%%%%%%%%%%%%%%%%%%%%%%%%%%
%%%%%%%%%%%%%%%%%%%%%%%%%%%%%
%%%%%%%%%%%%%%%%%%%%%%%%%%%%%%
%%%%%%%%% begin doc


\begin{document}
Assigned problems:
Chapter 1: 2, 4, 6, 7, 12, 15, 18, 26, 28
\subsubsection*{Problem 2}
\begin{problem}
  Prove whether the cross product is associative or not. 
  \begin{proof}
    We want to show that the cross product is not associative via a counterexample.\gline
    Let $\vA = \pvec{0, 1, 0}$, $\vB = \pvec{0, 1, 0}$, and $\vC = \pvec{1, 0, 0}$. Thus, we have that:
    \begin{align*}
      (\vA \times \vB)\times \vC &= \vec 0 \times \vC\\
      &= \vec 0\\
      \vA \times (\vB \times \vC) &= \vA \times \pvec{0, 0, 1}\\
      &= \pvec{1, 0, 0}
    \end{align*}
    Thus, because there exists an example for which the statement does not hold, the statement is false. 
  \end{proof}
\end{problem}
\subsubsection{Problem 4}
Use the cross product to find the components of the unit vector $\norm$ perpindicular to the shaded plane.
\gimage{image.png}{4cm}{Texbook diagram 1.11}
\begin{problem}
  Use the cross product to find the components of the unit vector $\norm$ perpindicular to the shaded plane. \gline
  We want to find a vector perpindicular to the plane that the surface of the triangle rests on. We just pick two linearly independent vectors:
  \begin{align*}
    \vA &= \pvec{-1, 2, 0}\\
    \vB &= \pvec{1, 0, -3}
  \end{align*}
  Taking the cross product gives us:
  \begin{align*}
    \vA \times\vB = \pvec{-6, -3, -2}
  \end{align*}
  Flipping the sign and taking the magnitude gives:
  \begin{align*}
    -(\vA \times\vB) &= (\vB \times \vA)\\
    &=\pvec{6, 3, 2}\\\\
    |\vB\times \vA| &= \sqrt{6^2 + 3^2 + 2^2}\\
    &= 7\\\\
    \norm &= \pvec{\frac{6}7,\frac{3}7, \frac{2}7}
  \end{align*}
  Thus, we have our unit vector. 
\end{problem}
\subsubsection*{Problem 6}
Prove that:
\begin{align*}
  [\vA\times(\vB\times\vC)]+[\vB\times(\vC\times\vA)]+[\vC\times(\vA\times\vB)] = \vZ
\end{align*}
Under what conditions is the cross product associative. 
\begin{problem}
  Consider that $\vA\times(\vB\times\vC)=\vB(\vA\cdot\vC)-\vC(\vA\cdot\vB)$. Writing out all parts of the equation in this notation yields:
  \begin{align*}
    &\vB(\vA\cdot\vC)-\vC(\vA\cdot\vB) + \vC(\vB\cdot\vA)-\vA(\vB\cdot\vC) + \vA(\vC\cdot\vB)-\vB(\vC\cdot\vA)\\
    &=\;\;\;\vB(\vA\cdot\vC)-\vB(\vC\cdot\vA)\\
    &\quad-\vC(\vA\cdot\vB)+ \vC(\vB\cdot\vA)\\
    &\quad-\vA(\vB\cdot\vC) + \vA(\vC\cdot\vB)
  \end{align*}
  Because dot product is commutative, all like terms cancel. Thus, the equation simplifies to $\vZ$. \gline
  The cross product is then associative if all terms are scalar multiples of each other, or the vectors form a basis for their space. In other words, this occurs when $\vA \times (\vB \times \vC) = \vZ$. 
\end{problem}
\subsubsection*{Problem 7}
Find the separation vector $\ar$ from the source point $\pvec{2, 8, 7}$ to the field point $\pvec{4, 6, 8}$. Determine it's magnitude $(r)$, and construct the unit vector $\arhat$. 
\begin{problem}
  First, we find our separation vector:
  \begin{align*}
    \ar &= \pvec{4 - 2, 6 -8 , 8-7}\\
    &= \pvec{2, -2, 1}
  \end{align*}
  Then, we find the magnitude:
  \begin{align*}
    r&= \sqrt{2^2 + (-2)^2 + 1}\\
    &= \sqrt{9}\\
    &= 3
  \end{align*}
  Then, that gives us the normal:
  \begin{align*}
    \arhat = \pvec{\frac{2}3, \frac{-2}3, \frac{1}3}
  \end{align*}
\end{problem}
\subsection*{Problem 12}
The height of a certain hill (in feet) is given by:
\begin{align*}
  h(x,y) = 10(2xy - 3x^2 - 4y^2 - 18x + 28y +12)
\end{align*}
where $y$ is the distance (in miles) north, $x$ the distance easy of South hadley.
\begin{enumerate}
  \item Where is the top of the hill located?
  \begin{problem}
    To solve this, we will take the gradient of this function and find the maximum point by solving the Hessian.
    \begin{align*}
      \prt{}{x}h(x,y) &= 10(2y - 6x - 18)\\
      &= 20(y - 3x - 9)\\
      \prt{}yh(x,y)&= 10(2x - 8y + 28)\\
      &= 20(x - 4y + 14)
    \end{align*}
    Solving the Hessian gives us:
    \begin{align*}
      \text{Det}\begin{bmatrix}
        -6 & 0\\
        0&-4
      \end{bmatrix}=24
    \end{align*}
    Now, we can solve for each function equaling zero, which gives us the representative equations:
    \begin{align*}
      3x &= y - 9\\
      x &= 4y -14\\
    \end{align*}
    Thus, substituting in $x$ gives:
    \begin{align*}
      3(4y-14)&= y -9\\
      12y - 42 &= y - 9\\
      11y &= 33\\
      y &= 3\\
      x &= -2
    \end{align*}
    Thus, we get that our point is a maximum with the hessian because both $f_{xx}$ and $f_{yy}$ are less than zero.\gline{}
    Via our hessian, we know that this a maximum. 
  \end{problem}
  \item How high is the hill?
  \begin{problem}
    If we plug in our $x$ and $y$, we get that the hill's top height is at 72. 
  \end{problem}
  \item How steep is the slope (in feet per mile) at a point 1 mile north and 1 mile east of South Hadley? In what direction is the slope steepest, at that point?
  \begin{problem}
    If we plug in our $x$ and $y$ for that point, we get that via the gradient the directional derivative is:
    \begin{align*}
      \{-220, 220\}
    \end{align*}
    The unit vector for this point is just $\{-\frac{\sqrt 2}2, \frac{\sqrt 2}2\}$, our direction (given a total steepness of 311.1269).
  \end{problem}
\end{enumerate}
\subsection*{Problem 15}
Calculate the divergence of the following vector functions:
\begin{enumerate}
  \item $\vv_a =\pvec{x^2, 3xz^2, -2xz}$
  \begin{problem}
    \begin{align*}
      \nabla&\cdot \vv_a\\
      \delfull &\cdot \pvec{x^2, 3xz^2, -2xz}\\
      \delinput{x^2}{2xz^2}{-2xz}&= (2x + 0 + -2x)\\
      &= 0
    \end{align*}
  \end{problem}
  \item $\vv_b = \pvec{xy, 2yz, 3xz}$
  \begin{problem}
    \begin{align*}
      \nabla&\cdot \vv_a\\
      \delfull&\cdot\pvec{xy, 2yz, 3xz}\\
      \delinput{xy}{2yz}{3xz}&=(y + 2z + 3z)
    \end{align*}
  \end{problem}
  \pagebreak
  \item $\vv_c = \pvec{y^2, 2xy+z^2, 2yz}$
  \begin{problem}
    \begin{align*}
      \nabla&\cdot \vv_a\\
      \delfull&\cdot\pvec{y^2, 2xy+z^2, 2yz}\\
      \delinput{y^2}{2xy+z^2}{2yz}&=(0+2x+2y)
    \end{align*}
  \end{problem}
\end{enumerate}
\subsection*{Problem 18}
Calculate the curls of the vector function in Prob. 1.15.
\begin{enumerate}
  \item $\vv_a =\pvec{x^2, 3xz^2, -2xz}$
  \begin{problem}
    \begin{align*}
      \delmatrix{x^2}{3xz^2}{-2xz}&\Rightarrow \pvec{0 - 4xz, 0 +2z, 3z^2-0}\\
      &\Rightarrow \pvec{-6xz, 2z, 3z^2}
    \end{align*}
  \end{problem}
  \item $\vv_b = \pvec{xy, 2yz, 3xz}$
  \begin{problem}
    \begin{align*}
      \delmatrix{xy}{2yz}{3xz}&\Rightarrow \pvec{0 - 2y, 0 - 3z, 0 - x}\\
      &\Rightarrow \pvec{-2y, -3z, -x}
    \end{align*}
  \end{problem}
  \item $\vv_c = \pvec{y^2, 2xy+z^2, 2yz}$
  \begin{problem}
    \begin{align*}
      \delmatrix{y^2}{2xy+z^2}{2yz}&\Rightarrow \pvec{2z - 2z, 0 - 0, 2y - 2y}\\
      &\Rightarrow\pvec{0, 0, 0}
    \end{align*}
  \end{problem}
\end{enumerate}
\pagebreak
\subsection*{Problem 26}
Calculate the Laplacian of the following function:
\newcommand{\laplace}{\prt{^2}{x^2}+\prt{^2}{y^2}+\prt{^2}{z^2}}
\newcommand{\prts}[2]{\prt{^2#1}{#2^2}}
\begin{enumerate}
  \item $T_a = x^2 + 2xy + 3z + 4$
  \begin{problem}
    We just apply the Laplacian to this function:
    \begin{align*}
      \cpr{\laplace}&\Rightarrow x^2 + 2xy+3z+4\\
      \prts{T_a}{x} &= 2\\
      \prts{T_a}{y}&= 0\\
      \prts{T_a}{z}&= 0\\
      2 + 0 + 0 &= 2
    \end{align*}
  \end{problem}
  \item $T_b = \sin x\sin y\sin z$
  \begin{problem}
    \begin{align*}
      \cpr{\laplace}&\Rightarrow\sin x\sin y\sin z\\
      \prts{T_b}{x}&= \prt{}{x}\cpr{\cos x\sin y\sin z} = -\sin x \sin y \sin z\\
      \prts{T_b}{y}&= \prt{}{y}\cpr{\sin x \cos y \sin z} = -\sin x \sin y \sin z\\
      \prts{T_b}{z}&= \prt{}{z}\cpr{\sin z \sin y\cos z} = -\sin x\sin y\sin z\\
      +\cdots &= 3\sin x\sin y \sin z
    \end{align*}
  \end{problem}
  \item $T_c = e^{-5x}\sin 4y\cos 3z$
  \begin{problem}
    \begin{align*}
      \cpr{\laplace}&\Rightarrow e^{-5x}\sin 4y\cos 3z\\
      \prts{T_c}{x}&= \prt{}{x}\cpr{-5e^{-5x}\sin 4y\cos 3z}= 25e^{-5x}\sin 4y\cos 3z\\
      \prts{T_c}{y}&= -16e^{-5x}\sin 4y\cos 3z\\
      \prts{T_c}{z}&= -9e^{-5x}\sin 4y\cos 3z\\
      0&= (25 -16 - 9)(25e^{-5x}\sin 4y\cos 3z)
    \end{align*}
  \end{problem}
  \pagebreak
  \item $\vv = \pvec{x^2, 3xz^2, -2xz}$
  \begin{problem}
    \begin{align*}
      \cpr{\laplace}&\Rightarrow \vv = \pvec{x^2, 3xz^2, -2xz}\\
      \cpr{\laplace}{x^2}&= 2\\
      \cpr{\laplace}{3xz^2}&= 6x\\
      \cpr{\laplace}{-2xz}&= 0\\
      \pvec{2, 6x, 0}
    \end{align*}
  \end{problem}
\end{enumerate}
\subsection*{Problem 28}
Prove the curl of the gradient is always zero. 
\begin{proof}
  We want to show that the curl of a gradient is always zero. That is, that $\nabla \times (\nabla f(\vv)) = 0$. Consider that:
  \begin{align*}
    \nabla f(\vv)&= \pvec{\prt{f}{x}, \prt{f}{y}, \prt{f}{z}}
  \end{align*}
  Thus, we have that:
  \begin{align*}
    \text{Det}\delmatrix{\prt{f}{x}}{\prt{f}{y}}{\prt{f}{z}}&= \cdots\\
    \cdots =&\pvec{ \prt{f^2}{yz}-\prt{f^2}{zy}, \prt{f^2}{xz}-\prt{f^2}{zx}, \prt{f^2}{xy}-\prt{f^2}{yx}}
  \end{align*}
  Note that for any variable $x_1, x_2$, we have that $\prt{f^2}{x_1x_2} = \prt{f^2}{x_2x_1}$. Thus, all of our terms cross cancel, and we are left with the vector:
  \begin{align*}
    \pvec{0, 0, 0}.
  \end{align*}
\end{proof}\vspace{0.4cm}
We check our solution with the function in Problem 1.11. Consider the function $x^2y^3z^4$. Our gradient is given by:
\begin{align*}
  \pvec{2xy^3z^4, 3x^2y^2z^4, 4x^2y^3 z^3}
\end{align*}
Taking our cross product with respect to the del operator gives us:
\begin{align*}
  \text{Det}\delmatrix{2xy^3z^4}{3x^2y^2z^4}{4x^2y^3 z^3} &= \cdots\\
  \cdots = \langle12x^2y^2z^3 - 12x^2y^2z^3&, 8xy^3z^3 - 8xy^3z^3, 6xy^2z^3 - 6xy^2z^3\rangle\\
  &= \pvec{0,0,0}
\end{align*}
\end{document}